% part: intuitionistic-logic
% chapter: propositions-as-types
% section: normalization

\documentclass[../../../include/open-logic-section]{subfiles}

\begin{document}

\olfileid{int}{pty}{nor}

\olsection{Normalization}

In this section we prove that, via some reduction order, any deduction
can be reduced to a normal deduction, which is called the
\emph{normalization property}. We will make use of the
propositions-as-types correspondence: we show that every proof term
can be reduced to a normal form; normalization for natural deduction
!!{derivation}s then follows.

Firstly we define some functions that measure the complexity of terms.
The \emph{length}~$\len{!A}$ of !!a{formula}s is defined by
\begin{align*}
  \len{p} &= 0\\
  \len{\lfalse} &= 0\\
  \len{!A \land !B} &= \len{!A} + \len{!B} + 1\\
  \len{!A \lor !B} &= \len{!A} + \len{!B} + 1 \\
  \len{!A \lif !B} &= \len{!A} + \len{!B} + 1.
\end{align*}
The complexity of a redex~$M$ is measured by its \emph{cut
  rank}~$\cutrank{M}$:
\begin{align*}
  \cutrank{(\lambd[\typeof{x}{!A}][\typeof{N}{!B}])Q} &=
  \len{!A} + \len{!B} + 1 \\
  \cutrank{\ande{i}{\andi{\typeof{M}{!A}}{\typeof{N}{!B}}}} &=
  \len{!A} + \len{!B} + 1 \\
  \cutrank{\ore{\ori{i}{!A_{i}}{\typeof{M}{!A_i}}}
    {\typeof{x_1}{!A_1}}{\typeof{N_1}{!C}}
    {\typeof{x_2}{!A_2}}{\typeof{N_2}{!C}}} &=
  \len{!A_1} + \len{!A_2} + 1
\end{align*}
The complexity of a proof term is measured by the most complex
redex in it, and $0$ if it is normal:
\[
  \maxrank{M} = \max \{\cutrank{N} | N \text{ is a sub term of } M \text{
    and is redex}\}
\]

\begin{lem}\ollabel{lem:subst}
  If $\Subst{M}{\typeof{N}{!A}}{\typeof{x}{!A}}$ is a redex and $M
  \not\ident x$, then one of the following cases holds:
  \begin{enumerate}
  \item $M$ is itself a redex, or
  \item $M$ is of the form $\ande{i}{x}$, and $N$ is of the form
    $\andi{P_1}{P_2}$
  \item $M$ is of the form $\ore{x}{x_1}{P_1}{x_2}{P_2}$, and $N$ is
    of the form $\ori{i}{}{Q}$
  \item $M$ is of the form $x Q$, and $N$ is of the form
    $\lambd[x][P]$
  \end{enumerate}
  In the first case, $\cutrank{\Subst{M}{N}{x}} = \cutrank{M}$; in the
  other cases, $\cutrank{\Subst{M}{N}{x}} = \len{!A}$.
\end{lem}
\begin{proof}
  Proof by induction on~$M$.
  \begin{enumerate}
  \item If $M$ is a single variable $y$ and $y \not\ident x$, then
    $\Subst{y}{N}{x}$ is $y$, hence not a redex.
  \item If $M$ is of the form $\andi{N_1}{N_2}$, or $\lambd[x][N]$, or
    $\ori{i}{!A}{N}$, then $\Subst{M}{\typeof{N}{!A}}{\typeof{x}{!A}}$ is
    also of that form, and so is not a redex.
  \item If $M$ is of the form $\ande{i}{P}$, we consider two cases.
    \begin{enumerate}
      \item If $P$ is of the form $\andi{P_1}{P_2}$, then $M \ident
        \ande{i}{\andi{P_1}{P_2}}$ is a redex, and clearly
        \[
        \Subst{M}{N}{x} \ident
        \ande{i}{\andi{\Subst{P_1}{N}{x}}{\Subst{P_2}{N}{x}}}
        \]
        is also a redex. The cut ranks are equal.
      \item If $P$ is a single variable, it must be $x$ to make the
        substitution a redex, and $N$ must be of the form
        $\andi{P_1}{P_2}$. Now consider
        \[
        \Subst{M}{N}{x} \ident \Subst{\ande{i}{x}}{\andi{P_1}{P_2}}{x},
        \]
        which is $\ande{i}{\andi{P_1}{P_2}}$. Its cut rank is equal to
        $\len{!A}$.
    \end{enumerate}
  \end{enumerate}
  The cases of $\ore{N}{x_1}{N_1}{x_2}{N_2}$ and $P Q$ are similar.
\end{proof}

\begin{lem}
  If $M$ contracts to $M'$, and $\cutrank{M} > \cutrank{N}$ for all
  proper redex sub-terms $N$ of $M$, then $\cutrank{M} >
  \maxrank{M'}$.
\end{lem}

\begin{proof}
  Proof by cases.
  \begin{enumerate}
  \item If $M$ is of the form $\ande{i}{\andi{M_1}{M_2}}$, then $M'$ is
    $M_i$; since any sub-term of $M_i$ is also proper
    sub-term of $M$, the claim holds.
  \item If $M$ is of the form
    $(\lambd[\typeof{x}{!A}][N])\typeof{Q}{!A}$, then $M'$ is
    $\Subst{N}{\typeof{Q}{!A}}{\typeof{x}{!A}}$.  Consider a redex in
    $M'$.  Either there is corresponding redex in $N$ with equal cut
    rank, which is less than $\cutrank{M}$ by assumption, or the cut
    rank equals $\len{!A}$, which by definition is less than
    $\cutrank{(\lambd[x^{!A}][N])Q}$.
  \item If $M$ is of the form
    \[
    \ore{\ori{i}{}{\typeof{N}{!A_i}}}
        {\typeof{x_1}{!A_1}}{\typeof{N_1}{!C}}
        {\typeof{x_2}{!A_2}}{\typeof{N_2}{!C}},
    \]
    then $M' \ident \Subst{N_i}{N}{\typeof{x_i}{!A_i}}$.  Consider a
    redex in~$M'$. Either there is corresponding redex in $N_i$ with
    equal cut rank, which is less than $\cutrank{M}$ by assumption; or
    the cut rank equals $\len{!A_i}$, which by definition is less than
    $\cutrank{\ore{\ori{i}{}{\typeof{N}{!A_i}}}
      {\typeof{x_1}{!A_1}}{\typeof{N_1}{!C}}
      {\typeof{x_2}{!A_2}}{\typeof{N_2}{!C}}}$.
  \end{enumerate}
\end{proof}

\begin{thm}
  All well-typed proof terms reduce to normal form; all !!{derivation}s reduce to
  normal !!{derivation}s.
\end{thm}

\begin{proof}
  For well-typed terms use the established result in
  \href{https://doi.org/10.2201/NiiPi.2005.2.4}{Tatsuta (2005),
  \emph{Second order permutative conversions with Prawitz's strong validity}},
  Definitions 2.3--2.4 and Theorem 4.15. Its language includes bottom,
  without adding bottom reduction rules.
  \begin{enumerate}

  \item Translate the well-typed term $M$ into that system:
  retain variable and lambda type annotations and use corresponding
  application, pair, projection, injection, case and bottom-elimination
  constructors. Rename bound variables freshly. Every typing rule and
  every beta or permutation step defined here is a typing rule or a
  reduction step there, also under compatible subterm closure.
  \item An infinite path here would therefore give an infinite path
  there, contradicting the cited strong normalization theorem.
  \begin{enumerate}
  \item Continuing reduction must reach a term with no available
  reduction, which is a normal form.

  \item Full normalization of natural deduction !!{derivation}s is
  independently established in
  \olref[pt][nor][thm]{thm:normalization}.
  The beta-redex rank calculation alone does not establish all
  permutation normalization.
  \end{enumerate}
  \end{enumerate}
\end{proof}

The fact that typed $\lambd$-terms normalize means that every such
term \emph{has a normal form}. If we think of $\beta$-reducing
$\lambd$-terms to their normal form as executing a computation and the
normal form as the value of the computation, this means that every
computation in the typed $\lambd$-calculus eventually halts an
produces a value. Moreover, type preservation ensures that the value
has the right type. E.g., if $N$ has type $!A \lif !B$ (i.e., it is a
function that is supposed to work for arguments of type~$!A$ and
produces values of type~$!B$) and we apply it to a term~$M$ of type~$!A$
(the input), $(NM)$ reduces to a term $N'$ (the output) and that
output is guaranteed to be of type~$!B$.

In fact, terms in the typed $\lambd$-calculus don't just reduce to
normal forms for one way of applying $\beta$-reductions to subterms,
but \emph{any} way of applying $\beta$-reductions to subtems
eventually terminates in a normal form. This property is called
\emph{strong normalization}. It means that not only is there a way to
guarantee that a computation eventually halts and produces a value,
but \emph{any} way to execute the computation eventually produces a
value.

How do we know that different ways of executing a computation produce
the \emph{same} value? So far we just know (from type preservation)
that all values it produces are of the same type. The typed
$\lambd$-calculus has another important property: it is
\emph{confluent} or \emph{Church-Rosser}. If $N \red N_1$ and $N \red
N_2$ then there is some term~$N'$ such that $N_1 \red N'$ and $N_2 \red
N'$. Recall that $N \red N'$ means: $N'$ results from zero or more
$\redone$-reduction steps. If $N_1$ is in normal form then the only
term $N'$ such that $N_1 \red N'$ is $N_1$ itself (via zero
$\redone$-reduction steps). It follows that if $N_1$ and $N_2$ are
both in normal form, then $N_1 = N' = N_2$. In other words, since
$\redone$ is strongly normalizing, every way of reducing a term
eventually results in a normal form. Since $\red$ is Church-Rosser,
any two normal forms are equal. Hence, computation in the typed
$\lambd$-calculus always terminates and the value (normal form) is the
same regardless of how the computation proceeds.

Proving that a relation is confluent is typically difficult. However,
we can easily establish the following weaker condition:

\begin{prop}[Weak Church-Rosser Property]
If $N \redone N_1$ and $N \redone N_2$ then there is a term $N'$ such
that $N_1 \red N'$ and $N_2 \red N'$.
\end{prop}

\begin{proof}
  The two terms $N_1$ and $N_2$ result from $N$ by replacing the redex
  $M_1$ or $M_2$ by its reductum, respectively. The redexes $M_1$ and
  $M_2$ are subterms of $N$, and so cannot overlap. This means that
  either one is a subterm of the other, or they occur in different
  non-overlapping places in $N$. Consider the latter case: $N$ is of
  the form $N(M_1,M_2)$. Replacing $M_1$ by its reductum $M_1'$
  results in $N_1$ which is $N(M_1',M_2)$, and replacing $M_2$ by its
  reductum $M_2'$ results in $N_2$ which is $N(M_1,M_2')$. Both reduce
  to $N(M_1',M_2')$. For the other case, suppose $M_2$ is a sub-term
  of~$M_1$ (the other case is symmetric) and distinguish cases
  according to what kind of redex $M_1$ is. E.g., if $M_1 \ident
  \proj{1}{\pair{O_1}{O_2}}$ it reduces to $O_1$. If $M_2$ is a
  subterm of $O_1$, then converting $M_2$ results in $O_1'$. So
  converting $M_1$ first and then $M_2$ results in $O_1'$. But
  converting $M_2$ first results in $\proj{1}{\pair{O_1'}{O_2}}$ and
  this reduces also to~$O_1'$.
\end{proof}

\begin{lem}[Newman's Lemma]
If $\redone$ is strongly normalizing and has the weak Church-Rosser
property, its closure $\red$ is confluent.
\end{lem}

\begin{proof}
  Since $\redone$ is strongly normalizing, any reduction sequence results
  in a normal form. Normal forms are unique iff $\red$ is confluent.
  So suppose some term $N$ has two distinct normal forms $N'$ and
  $N''$, the result of reduction sequences
  \begin{align*}
  N & = N_1' \redone N_2' \redone \dots \redone N_k' = N' \text{ and}\\
  N & = N_1'' \redone N_2'' \redone \dots \redone N_l'' = N''
  \end{align*}
  (The reduction sequence must have length $>0$ since otherwise $N$
  would already be in normal form.)

  If $N_2' = N_2''$ then it has two distinct normal forms ($N'$ and
  $N''$). If $N_2' \neq N_2''$ then by weak Church-Rosser there is
  some $N'''$ such that $N_2' \red N'''$ and $N_2'' \red N'''$. There is a normal form $Q$ of $N'''$. Since
  $N' \neq N''$, either $Q \neq N'$ or $Q \neq N''$.
  Consequently, either $N_2'$ or $N_2''$ has two distinct normal
  forms. We've shown that if $N$ has two distinct normal forms, there
  is an $N_1$ such that $N \redone N_1$ and $N_1$ has two distinct
  normal forms. Obviously we can keep going and produce $N \redone N_1
  \redone N_2 \redone \dots$, but this is impossible because $\redone$
  is strongly normalizing.
\end{proof}


\end{document}
